\section{Mixture proposals with the original model correction}
\label{sec:model-mixtures}

There is a continuous mixture already present in a particle filter whenever
the physical transition has a continuous density. Let
$\{x_{t-1}^j,W_{t-1}^j\}_{j=1}^N$ approximate the filtering distribution
after observation $t-1$, with $\sum_jW_{t-1}^j=1$. Propagating that measure
through the actual transition density gives
\begin{equation}
 a_t(x)=\sum_{j=1}^N W_{t-1}^j f_\theta(x\mid x_{t-1}^j),\qquad
 u_t(x)=g_\theta(y_t\mid x)a_t(x).
 \label{eq:physical-predictive-mixture}
\end{equation}
The first density predicts the state; the second incorporates the current
observation but is not normalized. The transition's process covariance is
part of the model. It is not an artificial bandwidth added to make an atom
cloud differentiable. When $f_\theta$ is Gaussian, $a_t$ is already a Gaussian
mixture, so Younis-style mixture calculus can be studied on this physical
predictive distribution directly.

Choose an evaluable proposal $q_t$ whose support contains that of $u_t$, and
draw $x_t^i\sim q_t$. The correct importance calculation is
\begin{equation}
 v_t^i=\frac{u_t(x_t^i)}{q_t(x_t^i)},\qquad
 \widehat z_t=\frac1N\sum_i v_t^i,\qquad
 W_t^i=\frac{v_t^i}{\sum_k v_t^k}.
 \label{eq:model-marginal-weights}
\end{equation}
Conditionally on the incoming cloud and a proposal fixed before these draws,
\begin{equation}
 \mathbb E_{q_t}[\widehat z_t\mid\mathcal F_{t-1}]
 =\int \frac{u_t(x)}{q_t(x)}q_t(x)\,\dd x
 =\int u_t(x)\,\dd x.
 \label{eq:conditional-model-normalizer}
\end{equation}
Here $\mathcal F_{t-1}$ includes the previous weighted particles and any
proposal design fixed from them and the observed data. A kernel can change
where particles are drawn without changing the numerator in this identity.
The cancellation fails if the denominator is merely the density of the
selected component rather than the actual marginal proposal, or if the
proposal is refitted to the same draws without a justified adaptive weighting
scheme.

This is the useful connection to the marginal particle filter.
\citet[Algorithm 2, Theorem 2, equations (9)--(12)]{lai2022vmpf} integrate out
the ancestor index and evaluate full transition and proposal mixtures. Their
estimator-coupling argument establishes an unbiased likelihood normalizer for
the specified recursion. It also explains a local Rao--Blackwell variance
reduction; the paper explicitly does not claim that the complete MPF always
has lower variance than every corresponding SMC run. Its variational
objective, equation (13), is $\mathbb E\log\widehat Z$, which is not
$\log Z$. Unbiased differentiation of that objective would not settle the
model-score question in Section~\ref{sec:question}.

The associated author implementation evaluates the transition mixture with
\texttt{reduce\_logsumexp}, adds the observation log density, and subtracts
the proposal log density \citep{laicode2022}. These operations support
\eqref{eq:model-marginal-weights}. They do not establish compatibility with
BayesFilter's analytical derivative, initialization, GPU/XLA, or tuning
requirements. The proposed generative mixture route is not implemented by
merely calling the existing observation-KDM endpoint.

The distinction between numerator and proposal also determines how much of
the current LEDH machinery can be reused. Its flow, UKF information, GenUT
moments, and OT couplings may supply proposal centers, covariances, or
deterministic transformations. The numerator must still use the intended
incoming filtering weights and the actual transition and observation
densities. Replacing those weights by uniform reset weights and those states
by uncorrected reset locations changes $a_t$. Equation~\eqref{eq:conditional-model-normalizer}
then certifies an integral of that changed approximation, not the original
model recursion. A deterministic transformation may be used as a proposal
only when its induced density, including every necessary Jacobian, is known.

A defensive mixture can protect proposal support when a guided component
misses a plausible mode. For a fixed $0<\eta<1$, let
\begin{equation}
 q_t(x)=(1-\eta)q_t^{\rm guided}(x)+\eta a_t(x).
 \label{eq:defensive-mixture}
\end{equation}
Since $q_t\ge\eta a_t$, the incremental weight satisfies
$v_t(x)\le g_\theta(y_t\mid x)/\eta$. A bounded observation likelihood
therefore gives a bounded incremental weight. It does not bound every score
derivative, and the mixing fraction remains a scope-specific proposal choice
to calibrate. An arbitrary positive fraction is not justified as a default.

\subsection{An analytically adapted Gaussian example}

The first implementation can be checked in a model where the best local
proposal is available exactly. Suppose
$f_\theta(x\mid x_{t-1}^j)=\mathcal N(x;m_j,Q)$ and
$g_\theta(y_t\mid x)=\mathcal N(y_t;Hx,R)$, with $Q$ and $R$ positive
definite. In the matrix-LGSSM test, $m_j=A(\theta)x_{t-1}^j$.
Multiplying the two Gaussian densities gives
\begin{equation}
 f_\theta(x\mid x_{t-1}^j)g_\theta(y_t\mid x)
 =c_j\,\mathcal N(x;\widetilde m_j,\widetilde Q),
 \label{eq:adapted-component-product}
\end{equation}
where completing the square, or conditioning a joint Gaussian, yields
\begin{align}
 V&=HQH^\mathsf T+R,& c_j&=\mathcal N(y_t;Hm_j,V),\\
 K&=QH^\mathsf T V^{-1},&
 \widetilde m_j&=m_j+K(y_t-Hm_j),\\
 \widetilde Q&=Q-QH^\mathsf T V^{-1}HQ.&
 \label{eq:adapted-component-parameters}
\end{align}
The component evidence $c_j$ tells us how plausible ancestor $j$ is after the
new observation. Its conditional Gaussian moves the proposal toward that
observation. Thus the normalized predictive update is the mixture
\begin{equation}
 q_t^*(x)=\sum_j\frac{W_{t-1}^jc_j}{z_t^*}
                   \mathcal N(x;\widetilde m_j,\widetilde Q),
 \qquad z_t^*=\sum_j W_{t-1}^jc_j.
 \label{eq:fully-adapted-mixture}
\end{equation}
Substituting $q_t^*=u_t/z_t^*$ into
\eqref{eq:model-marginal-weights} makes every incremental weight equal to
$z_t^*$. The current-step proposal variance vanishes conditionally on the
incoming cloud. Earlier particle approximation remains, so this is not an
exact finite-particle Kalman filter. The example isolates the improvement
available from conditional integration and proposal adaptation before any
extra kernel smoothing is introduced. In particular, $Q$ here is the physical
process covariance, not the carried UKF mark or an added kernel bandwidth.

\subsection{Retained ancestry and a flow proposal}

The same model correction can be formulated on a joint ancestor--state
space. This is useful when the flow depends on the selected ancestor and
its inverse cannot be marginalized cheaply. The distinction determines the
correct numerator as well as the denominator; the following existing
\texttt{MODEL-IS} derivation is retained for that reason.
